\documentclass[unicode,10pt,a4paper,oneside,numbers=endperiod,openany]{scrartcl}

\usepackage{ifthen}
\usepackage[utf8]{inputenc}
\usepackage{graphics}
\usepackage{graphicx}
\usepackage{textcomp}
\usepackage{tcolorbox}
\usepackage{hyperref}
\usepackage{svg}
\usepackage[section]{placeins}
\usepackage[backend=biber,style=numeric]{biblatex}
\addbibresource{sources.bib}

\pagestyle{plain}
\voffset -5mm
\oddsidemargin  0mm
\evensidemargin -11mm
\marginparwidth 2cm
\marginparsep 0pt
\topmargin 0mm
\headheight 0pt
\headsep 0pt
\topskip 0pt
\textheight 255mm
\textwidth 165mm

\newcommand{\duedate} {}
\newcommand{\setduedate}[1]{%
\renewcommand\duedate {Due date:~ #1}}
\newcommand\isassignment {false}
\newcommand{\setassignment}{\renewcommand\isassignment {true}}
\newcommand{\ifassignment}[1]{\ifthenelse{\boolean{\isassignment}}{#1}{}}
\newcommand{\ifnotassignment}[1]{\ifthenelse{\boolean{\isassignment}}{}{#1}}

\newcommand{\assignmentpolicy}{
\begin{tcolorbox}[colback=gray!10, colframe=black,
                  title={HPC Lab for CSE 2025 --- Submission Instructions},
                  fonttitle=\bfseries,
                  sharp corners,
                  boxrule=0.8pt,
                  left=5pt, right=5pt, top=5pt, bottom=5pt]
% \small
\textbf{Following these instructions is mandatory}: Submissions that do not
comply \textbf{will not be considered}.
\begin{itemize}
\item \textbf{Submission Platform}:
      Assignments must be submitted via
      \href{https://sam-up.math.ethz.ch/?lecture=401-3670-00}{SAM-UP}.
\item \textbf{Required Files}: Your submission must include:
      \begin{itemize}
      \item \textbf{Report}: A PDF summarizing your results and observations,
            based on this \LaTeX\ template.
            \begin{itemize}
            \item File name:
                  \texttt{project\_number\_lastname\_firstname.pdf}
            \item You are allowed to discuss all questions with anyone you like;
                  however: (i) your submission must list anyone you discussed
                  problems with and (ii) you must write up your submission
                  independently.
            \end{itemize}
      \item \textbf{Source Code Archive}: A \textbf{compressed tar archive},
            named \texttt{project\_number\_lastname\_firstname.tar.gz},
            containing:
            \begin{itemize}
            \item All source code files.
            \item Build files and scripts (e.g., \texttt{Makefile}, batch job
                  scripts, plot scripts, etc).
                  If you modified any provided files, ensure they still work as
                  expected.
            \item The full \LaTeX\ source of your report including all needed
                  files to build the report (figures, listings, etc).
            \end{itemize}
      \item \textbf{Size Limit}: The total submission (PDF + archive) must not
            exceed \textbf{25 MB}.
      \item \textbf{File Naming Rules}: Use only \texttt{ASCII} characters
            (avoid spaces, special symbols, and non-English letters).
      \end{itemize}
\item \textbf{Grading Process}: The TAs will evaluate your project based on:
      \begin{itemize}
      \item Your report write-up.
      \item Your code implementation.
      \item Benchmarking your code's performance.
      \end{itemize}
\end{itemize}
\end{tcolorbox}
}
\newcommand{\punkte}[1]{\hspace{1ex}\emph{\mdseries\hfill(#1~\ifcase#1{Points}\or{Points}\else{Points}\fi)}}


\newcommand\serieheader[6]{
\thispagestyle{empty}%
\begin{flushleft}
\includegraphics[width=0.4\textwidth]{ETHlogo_13}
\end{flushleft}
  \noindent%
  {\large\ignorespaces{\textbf{#1}}\hspace{\fill}\ignorespaces{ \textbf{#2}}}\\ \\%
  {\large\ignorespaces #3 \hspace{\fill}\ignorespaces #4}\\
  \noindent%
  \bigskip
  \hrule\par\bigskip\noindent%
  \bigskip {\ignorespaces {\Large{\textbf{#5}}}
  \hspace{\fill}\ignorespaces \large \ifthenelse{\boolean{\isassignment}}{\duedate}{#6}}
  \hrule\par\bigskip\noindent%  \linebreak
 }

\makeatletter
\def\enumerateMod{\ifnum \@enumdepth >3 \@toodeep\else
      \advance\@enumdepth \@ne
      \edef\@enumctr{enum\romannumeral\the\@enumdepth}\list
      {\csname label\@enumctr\endcsname}{\usecounter
        {\@enumctr}%%%? the following differs from "enumerate"
        \topsep0pt%
        \partopsep0pt%
        \itemsep0pt%
        \def\makelabel##1{\hss\llap{##1}}}\fi}
\let\endenumerateMod =\endlist
\makeatother


\usepackage{listings}
\usepackage{xcolor}

\definecolor{codegreen}{rgb}{0,0.6,0}
\definecolor{codegray}{rgb}{0.5,0.5,0.5}
\definecolor{codepurple}{rgb}{0.58,0,0.82}
\definecolor{backcolour}{rgb}{0.95,0.95,0.95}

\lstdefinestyle{mystyle}{
    backgroundcolor=\color{backcolour},
    commentstyle=\color{codegreen},
    keywordstyle=\color{magenta},
    numberstyle=\tiny\color{codegray},
    stringstyle=\color{codepurple},
    basicstyle=\ttfamily\footnotesize,
    breakatwhitespace=false,
    breaklines=true,
    captionpos=b,
    keepspaces=true,
    numbers=left,
    numbersep=5pt,
    showspaces=false,
    showstringspaces=false,
    showtabs=false,
    tabsize=2
}

\lstset{style=mystyle}

\usepackage{amsmath}
\usepackage{adjustbox}
\usepackage{listings}
\usepackage{xcolor}

\lstset{basicstyle=\ttfamily\footnotesize, breaklines=true}

\begin{document}

\setassignment
\setduedate{09 March 2026, 18:00~(Europe/Zurich)}

\serieheader{High-Performance Computing Lab for CSE}{2026}
            {Student: Benjamin Diethelm}
            {Discussed with: Lukas Ammann \linebreak \null\hfill Raphael Caixeta \linebreak \null\hfill Gian Laager}
            {Report for Project 1}{}
\newline



\section{Euler warm-up [10 points]}

\begin{enumerate}
\item This tool is used to dynamically load packages as required, minimizing version conflicts and ensuring only the necessary components are loaded. It also sets environment variables as needed. Modules can be loaded or unloaded, for example: \begin{lstlisting}
    module load stack/2025-06
    module unload stack/2025-06
\end{lstlisting}
 or any other module available on the system. Additionally, the commands \begin{lstlisting}
    module avail
    module spider <module name>
\end{lstlisting}
 can be used to display a list of available modules or to get more information on a specific module respectively.


\item This tool is employed to manage jobs on Euler or other supercomputers. Slurm is accessed via the login nodes and facilitates access to the compute nodes. Users can request specific hardware resources and submit code for execution. Slurm verifies user permissions for the requested hardware and queues the job. Once the hardware becomes available, the code is executed on the compute node.


\item The C++ code in \texttt{01/hello\_world.cpp} outputs \texttt{Hello World from \{\$HOSTNAME\}}, where \texttt{\{\$HOSTNAME\}} represents the system's hostname.


\item The script is located in \texttt{01/job\_1\_4.sh}, with its output and error logs in \texttt{01/1\_4.out} and \texttt{01/1\_4.err}.


\item The script is located in \texttt{01/job\_1\_5.sh}, with its output and error logs in \texttt{01/1\_5.out} and \texttt{01/1\_5.err}.
\end{enumerate}

\section{Performance characteristics [50 points]}

\subsection{Peak performance [16 Points]}

The AMD website (~\cite{amd_epyc_7763_spec},~\cite{amd_epyc_7H12_spec}) provides specifications for this model, while the Euler CPU Node Documentation~\cite{euler_docs} offers details specific to Euler VII. Reference~\cite{uops} is used for FMA and superscalability factors. The base clock frequency is taken from~\cite{amd_epyc_7H12_spec} and~\cite{amd_epyc_7763_spec}.

\begin{itemize}
\item Epyc 7H12 / Phase 1: \[P_{\mathrm{\text{Core}}} = 2 \times 4 \times 2 \times 2.6\mathrm{\text{GHz}} = 41.6\,\text{ GFlop/s }\] \[P_{\mathrm{\text{CPU}}} = 64 \times P_{\mathrm{\text{core}}} = 2662.4\,\text{ GFlop/s }\] \[P_{\mathrm{\text{node}}} = 2 \times P_{\mathrm{\text{CPU}}} = 5324.8\,\text{ GFlop/s }\] \[P_{\mathrm{\text{cluster}}} = 292 \times P_{\mathrm{\text{node}}} = 1554.8\,\text{ TFlop/s }\]


\item Epyc 7763 / Phase 2: \[P_{\mathrm{\text{Core}}} = 2 \times 4 \times 2 \times 2.45\mathrm{\text{GHz}} = 39.2\,\text{ GFlop/s }\] \[P_{\mathrm{\text{CPU}}} = 64 \times P_{\mathrm{\text{core}}} = 2508.8\,\text{ GFlop/s }\] \[P_{\mathrm{\text{node}}} = 2 \times P_{\mathrm{\text{CPU}}} = 5017.6\,\text{ GFlop/s }\] \[P_{\mathrm{\text{cluster}}} = 248 \times P_{\mathrm{\text{node}}} = 1244.3\,\text{ TFlop/s }\]


\item Total: \[1554.8\,\text{ TFlop/s } + 1244.3\,\text{ TFlop/s } = 2799.1\,\text{ TFlop/s }\]
\end{itemize}

\subsection{Memory Hierarchies [8 Points]}

\subsubsection{Cache and main memory size}

Using the commands provided in the task description, we obtain the data in Tab. \ref{tab:perf_mem_7763}. The output files are available as \texttt{02/Memory\_7H12.txt} and \texttt{02/Memory\_7763.txt}, while the corresponding topology plots are included as Fig. \ref{fig:epyc_7h12} and Fig. \ref{fig:epyc_7763}. The job files used to produce this output are \texttt{02/job\_2\_1.sh} and \texttt{02/job\_2\_2.sh}.

The key difference is that the EPYC 7H12 shares L3 cache across 4 cores while the EPYC 7763 shares it across 8 cores. Main memory is partitioned into 4 NUMA nodes per CPU, yielding 8 NUMA nodes per node, for both phase 1 and 2.

\begin{table}[t]
    \centering
    \maxsizebox{\textwidth}{!}{\includegraphics{/home/benjamin/code/ttt/submission/generated/15853960861298318217.pdf}}
    \caption{Memory hierarchy of an Euler VII (phase 1/2) node with an AMD EPYC 7H12/7763 CPU. Each core has its own L1 and L2 cache. L3 cache is shared among groups of 4 cores in phase 1 and in groups of 8 cores in phase 2.\label{tab:perf_mem_7763}}
\end{table}
 

\begin{figure}[t]
    \centering
    \maxsizebox{\textwidth}{!}{\includegraphics{/home/benjamin/code/ttt/submission/generated/6377526862130354073.pdf}}
    \caption{Topology of an Euler VII (phase 1) node with an AMD EPYC 7H12 CPU.\label{fig:epyc_7h12}}
\end{figure}
 

\begin{figure}[t]
    \centering
    \maxsizebox{\textwidth}{!}{\includegraphics{/home/benjamin/code/ttt/submission/generated/16156204331478293454.pdf}}
    \caption{Topology of an Euler VII (phase 2) node with an AMD EPYC 7763 CPU.\label{fig:epyc_7763}}
\end{figure}
 

\subsection{Bandwidth: STREAM benchmark [10 Points]}

The following resources are available: source code (\texttt{02/stream.c}), batch scripts (\texttt{02/job\_2\_3.sh}, \texttt{02/job\_2\_4.sh}), output files (\texttt{02/2\_3\_7H12.out}, \texttt{02/2\_3\_7763.out}).

Tab. \ref{tab:stream_7h12} and Tab. \ref{tab:stream_7763} show the output from STREAM. Across all measurements, phase 2 shows better results. For both phases, Add and Triad are roughly consistent, while Scale is slightly lower. Copy throughput is notably higher for both, at approximately \(31.0\,\text{ GB/s}\) for phase 1 and \(36.6\,\text{ GB/s}\) for phase 2.

\begin{table}[t]
    \centering
    \maxsizebox{\textwidth}{!}{\includegraphics{/home/benjamin/code/ttt/submission/generated/13995748083104300913.pdf}}
    \caption{STREAM benchmark results for AMD EPYC 7H12 (Euler VII Phase 1).\label{tab:stream_7h12}}
\end{table}
 

\begin{table}[t]
    \centering
    \maxsizebox{\textwidth}{!}{\includegraphics{/home/benjamin/code/ttt/submission/generated/5936117419463465897.pdf}}
    \caption{STREAM benchmark results for AMD EPYC 7763 (Euler VII Phase 2).\label{tab:stream_7763}}
\end{table}
 

\subsection{Performance model: A simple roofline model [16 Points]}

The plot was generated using \texttt{02/roofline\_plot.py}, refer to Fig. \ref{fig:roofline_plot}. The script also calculates the minimum Operational Intensity required to achieve compute-bound performance. For Euler VII phase 1, the value is \(1.34\,\frac{\text{ Flops}}{\text{Byte}}\), and for phase 2, it is \(1.07\,\frac{\text{ Flops}}{\text{Byte}}\).

\begin{figure}[t]
    \centering
    \maxsizebox{\textwidth}{!}{\includegraphics{/home/benjamin/code/ttt/submission/generated/15401066031855344164.pdf}}
    \caption{Roofline model plot for Euler VII nodes.\label{fig:roofline_plot}}
\end{figure}
 

\section{Auto-vectorization [10 points]}

\begin{enumerate}
\item Data alignment refers to storing objects at memory addresses that are multiples of their own size. Misaligned memory access can result in significant performance penalties. SIMD/SSE instructions prefer 16-byte alignment for optimal efficiency. Data alignment can be enforced using \texttt{\_\_declspec(align(n))} or \texttt{\_\_builtin\_assume\_aligned}, or it may be handled automatically by the compiler.


\item Possible reasons include: \begin{itemize}
\item Non-contiguous memory access


\item Data-dependencies between loop iterations


\item Pointer aliasing


\item Non-countable loops


\item Function calls within loops


\item Mixing data types
\end{itemize}


\item Possible ways include: \begin{itemize}
\item \texttt{\#pragma ivdep} (ignore assumed dependencies)


\item \texttt{\#pragma vector always} (force vectorization)


\item \texttt{\#pragma vector align} (assert the data is aligned)


\item \texttt{\#pragma loop count(n)} (provide typical trip count for loop)


\item \texttt{\#pragma omp simd} (force SIMD)
\end{itemize}


\item Strip-mining splits a loop into a vectorized part and a scalar cleanup loop. Loop blocking partitions the iteration space into smaller blocks to improve cache reuse. Loop interchange reorders nested loops to obtain stride-1 memory access patterns. Loop peeling performs a few scalar iterations until an aligned boundary is reached.


\item We can use user-mandated vectorization with \texttt{\#pragma omp simd}, which forces vectorization and issues a warning if it fails. For more control, SIMD-enabled functions like \texttt{\_\_attribute\_\_((vector))} or \texttt{\#pragma omp declare simd} allow writing scalar-style functions that the compiler can transform into short vector variants. Vector intrinsics, e.g. \texttt{\_mm\_add\_ps} or inline assembly give full manual control.
\end{enumerate}

\pagebreak

\section{Matrix multiplication optimization [30 points]}

\subsection{dgemm\_jki.c}\label{sec:jki} 

The loop order was modified to \texttt{jki} to enhance cache locality, following standard GEMM optimization principles~\cite{goto_gemm,blis_toms}. Fig. \ref{fig:variants} compares this approach to the naive implementation, OpenBLAS, and other optimized versions.

The performance of \texttt{dgemm\_jki.c} demonstrates significant improvement over the naive implementation, achieving \(13.21\,\text{ GFlop/s}\) compared to \(3.41\,\text{ GFlop/s}\) for the naive version. This 3.9x speedup results from the JKI loop order's stride-one access pattern for matrix \texttt{C} in the innermost loop, which dramatically reduces cache misses compared to the naive IJK ordering. Despite non-contiguous access for matrix \texttt{A}, the JKI ordering achieves 26.3\% of peak performance on average versus only 4.6\% for the naive implementation.

\begin{figure}[t]
    \centering
    \maxsizebox{\textwidth}{!}{\includegraphics{/home/benjamin/code/ttt/submission/generated/15520668701961763214.pdf}}
    \caption{Performance comparison of the naive, jki, OpenBLAS and 1/2/3-level blocking versions.\label{fig:variants}}
\end{figure}
 

\subsection{dgemm\_blocked.c}

The final optimized file is included as \texttt{dgemm\_blocked.c}.

\subsubsection{System}\label{sec:system} 

Unless otherwise specified, all tests were conducted on a single core of an AMD EPYC 7763. The GCC compiler (version 13.2.0) was used, with flags \texttt{-O3 -march=znver3 -ffast-math -funroll-loops}~\cite{gcc_manual}. The default optimization strategy employed was 2-level blocking, as described in Section \ref{sec:overview}.

\subsubsection{Overview}\label{sec:overview} 

The optimizations implemented are:

\begin{itemize}
\item \textbf{Multi-level blocking for cache hierarchy}: The code implements three configurable blocking strategies (1-level, 2-level, and 3-level blocking). The 1-level blocking targets the L3 cache with block size \texttt{S3}, 2-level blocking adds L2 cache optimization with block size \texttt{S2}, and 3-level blocking further incorporates L1 cache optimization with block size \texttt{S1}. This multi-level cache blocking approach is a standard technique in high-performance GEMM implementations~\cite{goto_gemm,blis_toms}.


\item \textbf{Hierarchical block size optimization}: A systematic tuning process was employed to find optimal block sizes for each cache level. This kind of empirical parameter search is a common strategy in high-performance BLAS implementations~\cite{atlas}.


\item \textbf{Data packing for spatial locality}: The micro-kernel packs both matrices A and B into contiguous local buffers with optimal memory layouts, eliminating strided memory access patterns and ensuring sequential access during computation~\cite{goto_gemm,blis_toms}.


\item \textbf{Register-level accumulation}: The micro-kernel loads the C tile into a local accumulator array before computation. All updates to C are performed on this local buffer, significantly reducing memory writes~\cite{goto_gemm,blis_toms}.


\item \textbf{Loop dependency annotation}: The code uses \texttt{\#pragma GCC ivdep} in all innermost loops to inform the compiler that loop iterations are independent and can be safely vectorized~\cite{gcc_manual}.


\item \textbf{Pointer aliasing control}: All function parameters use the \texttt{restrict} keyword to guarantee to the compiler that pointers do not alias, enabling more aggressive optimizations~\cite{gcc_manual}.


\item \textbf{Compiler flag optimization}: The script \texttt{04/flag\_compare.sh} compares different compiler optimization flags to quantify the individual impact of each flag~\cite{gcc_manual,atlas}.
\end{itemize}

\subsubsection{Analysis}\label{sec:analysis} 

This section analyzes the impact of the optimizations introduced in \texttt{dgemm\_blocked.c} and relates them to the overall performance comparison shown in Fig. \ref{fig:variants}. In all plots, two complementary metrics are reported: the average percentage of peak (a good indicator of sustained performance across problem sizes) and the peak throughput in \(\text{GFlop/s}\) (which highlights best-case behavior under favorable cache and alignment conditions). Overall, the structure of the optimized implementation (cache blocking + packing + micro-kernel accumulation) is aligned with widely used high-performance GEMM designs~\cite{goto_gemm,blis_toms}.

The incremental results in Fig. \ref{fig:steps} show that the largest single improvement comes from switching the loop order to \texttt{jki}. Subsequent steps such as packing and register accumulation improve the sustained average further, but their effect is smaller and less uniform across sizes.

\begin{figure}[t]
    \centering
    \maxsizebox{\textwidth}{!}{\includegraphics{/home/benjamin/code/ttt/submission/generated/11098391828362779878.pdf}}
    \caption{Performance comparison of the incremental optimizations in \texttt{dgemm\_blocked.c}.\label{fig:steps}}
\end{figure}
 

For 2-level blocking, Fig. \ref{fig:l2} exhibits a pronounced optimum at \(S_{2} = 64\) (with \(S_{3}\) fixed at the best value from Fig. \ref{fig:l3}). This sharp peak indicates that once the L2 block becomes too large, the packed panels and the active C tile no longer fit comfortably in L2; when it is too small, the overhead becomes dominant. The selected \(S_{2} = 64\) therefore represents a sweet spot between locality and overhead on the EPYC 7763 core.

For 3-level blocking, Fig. \ref{fig:l1l2_slices} and Fig. \ref{fig:l1l2_heatmap} confirm that only a narrow region of \(\left( S_{1},S_{2} \right)\) combinations performs well, with the best result at \(S_{1} = 32\) and \(S_{2} = 64\). As mentioned in Section \ref{sec:overview}, various compiler flags have been tested. Fig. \ref{fig:flags} suggests that \texttt{-O3 -march=znver3 -ffast-math -funroll-loops} yields the best result.

\begin{figure}[t]
    \centering
    \maxsizebox{\textwidth}{!}{\includegraphics{/home/benjamin/code/ttt/submission/generated/8757924707434660522.pdf}}
    \caption{Average percentage of peak performance as a function of \(S_{3}\) for 1-level blocking. Best result at \(S_{3} = 1152\) achieving 34.81\% of peak.\label{fig:l3}}
\end{figure}
 

\begin{figure}[t]
    \centering
    \maxsizebox{\textwidth}{!}{\includegraphics{/home/benjamin/code/ttt/submission/generated/13434767790748840726.pdf}}
    \caption{Average percentage of peak performance as a function of \(S_{2}\) for 2-level blocking, with \(S_{3} = 1152\) fixed. The optimal \(S_{2} = 64\) achieves 51.4\% of peak.\label{fig:l2}}
\end{figure}
 

\begin{figure}[t]
    \centering
    \maxsizebox{\textwidth}{!}{\includegraphics{/home/benjamin/code/ttt/submission/generated/9854709472188265678.pdf}}
    \caption{Average percentage of peak performance for 3-level blocking as a function of \(S_{1}\) and \(S_{2}\). Best result at \(S_{1} = 32\), \(S_{2} = 64\) achieving 25.26\% of peak.\label{fig:l1l2_slices}}
\end{figure}
 

\begin{figure}[t]
    \centering
    \maxsizebox{\textwidth}{!}{\includegraphics{/home/benjamin/code/ttt/submission/generated/2061434689125365536.pdf}}
    \caption{Heatmap of average percentage of peak performance for 3-level blocking as a function of \(S_{1}\) and \(S_{2}\).\label{fig:l1l2_heatmap}}
\end{figure}
 

\begin{figure}[t]
    \centering
    \maxsizebox{\textwidth}{!}{\includegraphics{/home/benjamin/code/ttt/submission/generated/9161148123824896444.pdf}}
    \caption{Performance comparison of various compilation flags for \texttt{dgemm\_blocked.c}.\label{fig:flags}}
\end{figure}
 

 It also sets environment variables as 

\printbibliography

\end{document}
